% vim: ai sw=2 spell spelllang=cs

\begin{frame}{Tracery}

  \begin{itemize}
    \item Sledují události a vypisují je
    \item \cmd{strace}: sleduje systémová volání
    \item \cmd{ltrace}: sleduje knihovní (příp.\ i systémová) volání

    \pause

    \item Používají systémové volání \code{ptrace}
    \begin{itemize}
      \item \cmd{ltrace} si nastaví breakpointy na všechny externí symboly z
	ELFu
      \item \cmd{strace} jednoduše použije \code{PTRACE_SYSCALL}
    \end{itemize}

    \pause

    \item Spuštění: \cmd{strace/ltrace binarka --volby}
    \item Další parametry
    \begin{itemize}
      \item \texttt{-e}: filtr (\cmd{\-e write}, \cmd{\-e trace=network})
      \item \texttt{-f}: sleduje i potomky
      \item \texttt{-o}: výstup do souboru
    \end{itemize}
  \end{itemize}

\end{frame}

\begin{frame}{Úkol}

  \heading{Práce s tracery}
  \begin{enumerate}
    \item Projděte si \iffimuni
      \file{pb173-bin/09/tracer.c}
    \else\ifpcdirlinux
      \file{04/main.c}
    \else
      \file{06/tracer.c}
    \fi\fi
    \item Přeložte

    \item Spusťte nejdříve s \cmd{ltrace}
    \begin{itemize}
      \item Kolik je volání \code{fwrite} do \file{libc}?
    \end{itemize}

    \item Spusťte s \cmd{strace}
    \begin{itemize}
      \item Kolik je volání \code{write} do jádra?
      \item $\Rightarrow$ \file{libc} bufferuje
    \end{itemize}

    \item Pomocí \cmd{ltrace} zjistěte, jak velký buffer \file{libc} používá
    \begin{itemize}
      \item Přidejte cyklus, který vypisujte znak pomocí \code{fwrite}
      \item Použijte parametr \cmd{\-S}
    \end{itemize}

    \item Pomocí \cmd{strace} ověřte totéž
    \begin{itemize}
      \item Pozorujte třetí parametr \code{write}
    \end{itemize}
  \end{enumerate}

\end{frame}

